\section*{Overview}

In contest geometry, line intersection usually means one of two things: deciding whether two segments intersect, or
computing the actual crossing point of two non-parallel lines. The first task should stay exact whenever possible; the
second often needs a carefully isolated floating-point step.

\section*{When to Use It}

Use these routines when:

\begin{itemize}
  \item the problem asks whether paths, walls, rays, or polygon edges cross,
  \item boundary contact counts as an intersection,
  \item you need a robust primitive before moving on to polygon or sweep-line logic.
\end{itemize}

\section*{Core Idea}

Two segments intersect if each segment straddles the line through the other, with collinear overlap handled
separately. Orientation tests are enough for the discrete predicate.

For the actual crossing point of two non-parallel infinite lines, solve the linear equations of the two lines. That
part can use floating-point arithmetic after the combinatorial case split is done.

\section*{Key Insight}

Most bugs come from treating all cases the same. Proper intersections, touching endpoints, and collinear overlap should
be separated mentally even if the final code merges them into one predicate.

\section*{Operations / Main Technique}

\begin{itemize}
  \item orientation of endpoint triples,
  \item on-segment checks for collinear cases,
  \item one predicate for segment intersection,
  \item one helper for the intersection point of non-parallel lines.
\end{itemize}

\paragraph{Worked example.}
If the endpoints of segment \(ab\) lie on opposite sides of line \(cd\), and the endpoints of \(cd\) lie on opposite
sides of line \(ab\), the segments cross properly. If one orientation is zero, it becomes a boundary-touch case.

\section*{Correctness Intuition}

Orientation determines the side of a line. For a proper intersection, each segment's endpoints must lie on different
sides of the other segment's supporting line. Collinear cases are not captured by side changes, so they are handled by
explicit range overlap.

\section*{Complexity Analysis}

All primitives here run in \(O(1)\) time and \(O(1)\) memory.

\section*{Implementation}

The reference code gives:

\begin{itemize}
  \item \texttt{segments\_intersect} with inclusive boundary handling,
  \item \texttt{line\_intersection} for two non-parallel infinite lines.
\end{itemize}

That is the split I use most often in contests: exact predicate first, coordinate computation second.

\section*{Common Pitfalls}

\begin{itemize}
  \item Forgetting that collinear overlap is still an intersection.
  \item Using a floating-point epsilon for problems that can be solved exactly with integer predicates.
  \item Not deciding whether touching at one endpoint should count.
  \item Computing the intersection point before checking for parallel lines.
\end{itemize}

\section*{Variants / Extensions}

\begin{itemize}
  \item Ray or half-line intersection.
  \item Segment-polygon intersection.
  \item Sweep-line event ordering, which still relies on the same primitives.
\end{itemize}

\section*{Practice Problems}

\begin{itemize}
  \item Determine whether two segments intersect.
  \item Count how many query segments intersect a polygon boundary.
  \item Compute the crossing point of two non-parallel lines.
\end{itemize}

\section*{References}

\begin{itemize}
  \item \href{https://cp-algorithms.com/geometry/segments-intersection.html}{cp-algorithms: Intersection of Segments}.
  \item \href{https://usaco.guide/plat/geo-prims}{USACO Guide: Geometry Primitives}.
  \item \href{https://codeforces.com/catalog}{Codeforces Catalog}.
\end{itemize}
